% Generated draft: requires mathematical and source review.
\documentclass[11pt]{article}
\usepackage[margin=25mm]{geometry}
\usepackage{fontspec}
\setmainfont{Latin Modern Roman}
\usepackage{amsmath,amssymb,mathtools,unicode-math}
\setmathfont{Latin Modern Math}
\usepackage{enumitem,needspace,xurl,hyperref}
\setlength{\parindent}{0pt}
\setlength{\parskip}{5pt}
\setlength{\emergencystretch}{3em}
\newcommand{\N}{\mathbb{N}}
\newcommand{\Z}{\mathbb{Z}}
\newcommand{\Q}{\mathbb{Q}}
\newcommand{\R}{\mathbb{R}}
\newcommand{\C}{\mathbb{C}}
\newcommand{\F}{\mathbb{F}}
\newcommand{\A}{\mathbb{A}}
\newcommand{\PP}{\mathbb{P}}
\newcommand{\E}{\mathbb{E}}
\newcommand{\eps}{\varepsilon}
\newcommand{\dd}{\,\mathrm{d}}
\begin{document}
\section*{UnsolvedMath ID: 20002813}
\textbf{Source problem code:} catalogue entry
\subsection{Definitions and mathematical statement}
Let $G$ be finite and let $(\mathcal X,\mathcal K)$ be a supercharacter theory: $\mathcal X$ partitions the complex irreducible characters, $\mathcal K$ partitions $G$ into superclasses, the numbers of blocks agree, and the canonical characters
\[
 \sigma_A=\sum_{\chi\in A}\chi(1)\chi\qquad(A\in\mathcal X)
\]
are constant on superclasses and form their class-function basis. The identity singleton and trivial-character conditions use the usual convention. The question asks when all pointwise products have nonnegative integer coefficients in this canonical basis; zero coefficients are allowed by ``positive integer combination.''

Write $d_\chi=\chi(1)$ and let $m_{\chi,\psi}^{\theta}=\langle\chi\psi,\theta\rangle_G$ be the tensor-product multiplicity. Define
\[
 D_C=\sum_{\theta\in C}d_\theta^2,\qquad
 N_{AB}^{C}=\sum_{\substack{\chi\in A,\ \psi\in B\\\theta\in C}}
          d_\chi d_\psi d_\theta\,m_{\chi,\psi}^{\theta}.
\]
The aid's exact formula is
\[
 \sigma_A\sigma_B=\sum_{C\in\mathcal X}a_{AB}^{C}\sigma_C,
 \qquad a_{AB}^{C}=N_{AB}^{C}/D_C.
\]
Superclass-function closure ensures this expansion exists, and character orthogonality gives the coefficient. Hence all coefficients are automatically nonnegative rational, and the requested property is exactly
\[
 D_C\mid N_{AB}^{C}\qquad(A,B,C\in\mathcal X).
\]

The aid reports integrality for every abelian supercharacter theory and every maximal theory. In the abelian case each degree is one, so the common coefficient on a character block is an integer tensor-product count. For the maximal theory of a nontrivial group, with $\sigma=\rho_{\mathrm{reg}}-1$,
\[
 \sigma^2=(|G|-1)1+(|G|-2)\sigma.
\]
For the minimal theory, however, the canonical singleton is $d_\chi\chi$, not the unscaled irreducible character. In $A_5$, a degree-three irreducible satisfies $\chi_3^2=1+\chi_3+\chi_5$, giving
\[
 \sigma_3^2=9\sigma_1+3\sigma_3+\tfrac95\sigma_5.
\]
Thus minimal-theory integrality is not automatic. The original remark asserts it for minimal theories of $S_n$, but that assertion conflicts with this canonical normalization: the standard degree-three character of $S_4$ has square $1+\chi_2+\chi_3+\chi'_3$, producing coefficient $9/2$ on $\sigma_2$. With unscaled irreducibles ordinary tensor multiplicities are integral; these normalizations must not be interchanged.

For another supercharacter normalization $\tau_A=c_A\sigma_A$, $c_A\ne0$, if $\tau_A\tau_B=\sum_C n_{AB}^{C}\tau_C$, then
\[
 a_{AB}^{C}=n_{AB}^{C}\frac{c_C}{c_Ac_B}.
\]
Consequently known integral products for orbit-normalized algebra-group supercharacters do not alone answer the canonical-$\sigma$ question; the scaling factors must pass the divisibility test. The original remark also hopes for an assertion about four supercharacter theories of $S_n$ and asks about natural algebra-group characters. It does not specify which theories or the relevant $n$ range, and the aid does not classify those remaining canonical cases.
\subsection{Short English statement}
Canonical supercharacter products have nonnegative rational coefficients, integral exactly under an explicit degree-weighted divisibility condition. Abelian and maximal theories pass, while minimal examples fail. Rescaling natural algebra-group characters can change integrality, so normalization is essential.
\subsection{Sources}
\begin{itemize}
\item American Institute of Mathematics workshop \emph{Supercharacters and combinatorial Hopf algebras}, Problem 1.1, \url{https://aimath.org/WWN/supercharacters/supercharacters.pdf}.
\item A. Hendrickson, \emph{Supercharacter theories of cyclic $p$-groups}, as named in the original remark; no year or publication details are supplied. The remark mentions maximal/minimal and four-theory symmetric-group questions and natural algebra-group characters.
\item Local UnsolvedMath registry, permanent ID 20002813, source code \texttt{catalogue entry}, canonical location \texttt{aim-representation-theory-notes.json}, \texttt{notes[84]}, record dated 14 August 2026. The original literature field is empty.
\item Supplied research attempt 1: partial result, problem status at run partially solved, rigor labelled proved, novelty confidence low. Its coefficient/divisibility and rescaling criterion, with the $A_5$ example, is machine-generated research aid, not peer reviewed; independently verify before citation. The symmetric-group normalization inconsistency is stated explicitly rather than silently retaining a false canonical claim.
\end{itemize}
\end{document}
